\documentclass[11pt, a4paper]{article}
% \documentclass{article}
\usepackage{tikz}
\usepackage{amsmath}
\usepackage[utf8]{inputenc}
\usepackage[hungarian]{babel}
\usetikzlibrary{calc}

\title{%
  Antwort von \\
  \large Janos Sarközi}


\begin{document}
\maketitle
\section{Vorbemerkungen}
Um aus der Gleichung 
\begin{equation}
    -\frac{q_0\cdot l^3}{24E\cdot I_y} + \frac{\mu_A\cdot l}{6E\cdot I_y} +
    \frac{\mu_B\cdot l}{3E\cdot I_y} = 0
\end{equation}
die Gleichung
\begin{equation}
    - \frac{q_0\cdot l}{8} + \frac{\mu_A}{2} + \mu_B = 0
\end{equation}
zu bekommen mache ich paar Vorbemerkungen. Dafür betrachten wir erst eine einfache Gleichung
\[
5x = 0
\]
und stellen die Frage: Welches $x$ löst diese Gleichung? Die Antowrt ist $x =
0$. Dabei ist wichtig zu wissen, dass die Zahl fünf ungleich Null ist. Jetzt
gehen wir ein Schritt weiter und betrachten die Gleichung
\[
5x + 5y = 0
\]
Hier muss für $x$ und $y$ die Gleichung $x+y=0$ gelten. Warum? Wir können die
obige Gleichung umschreiben:
\[
5(x + y) = 0
\]
Damit sind wir in der ähnlichen Situation wie bei der Gleichung $5x=0$. Jetzt
der letzte Schritt für die Vorbereitungen. Dafür machen wir eine Gleichung mit
drei Termen, wie bei deiner Aufgabe, nur nicht so mit krassen Zeichen wir bei
dir. Dabei sind die Buchstaben $a, b$ und $c$ beliebige Konstanten.
\[
    -\frac{a\cdot x}{b\cdot c} + \frac{a\cdot y}{b\cdot c} + \frac{a\cdot
    z}{b\cdot c} = 0
\]
\newpage
Alle drei Terme haben die Konstante $a$ im Zähler. Damit können wir diese
Konstante aus allen drei Thermen herausziehen.
\[
    a\left(-\frac{x}{b\cdot c} + \frac{y}{b\cdot c} + \frac{z}{b\cdot c}\right)
    = 0
\]
Nun entdecken wir was Ähnliches mit den Konstanten $b$ und $c$, die im Nenner
sind. Damit erhalten wir:
\[
    \frac{a}{b\cdot c}(-x + y + z) = 0
\]
Dies ist die Distributivität der reellen Zahlen. Normalerweise haben wir die
Gleichung in der Form
\[
    a(x+y) = ax+ay
\]
Wir haben also eher die einfachere Variante im Kopf: Wir lösen die Klammer auf.
Wenn wir die Gleichung aber umdrehen, dann führen wir die Klammer ein. Dies ist
meiner Meinung nach etwas schwieriger. Damit sind die Vorbereitungen
abgeschlossen.
\section{Antwort}
Um aus der Gleichung (1) die Gleichung (2) zu erhalten benutzen wir mehrmals die
Distributivität der reellen Zahlen. Ich rechne mal das in Zeitlupe vor:
\[
    -\frac{q_0\cdot l^3}{24E\cdot I_y} + \frac{\mu_A\cdot l}{6E\cdot I_y} +
    \frac{\mu_B\cdot l}{3E\cdot I_y} = 0
\]
Wir sehen in allen drei Termen die Variable $l$, dies können wir aus allen drei
Termen herausziehen:
\[
    l\left(-\frac{q_0\cdot l^2}{24E\cdot I_y} + \frac{\mu_A}{6E\cdot I_y} +
    \frac{\mu_B}{3E\cdot I_y}\right) = 0
\]
Jetzt entdecken wir das auch für $E$ und $I_y$. Damit haben wir
\[
    \frac{l}{E\cdot I_y}\left(-\frac{q_0\cdot l^2}{24} + \frac{\mu_A}{6} +
    \frac{\mu_B}{3}\right) = 0
\]
Nun haben wir noch glück und können die Zahl drei auch vor die Klammer Ziehen.
\[
    \frac{l}{3E\cdot I_y}\left(-\frac{q_0\cdot l^2}{8} + \frac{\mu_A}{2} +
    \mu_B\right) = 0
\]
Damit ich aus der Gleichung (1) die Gleichung (2) folgern kann, dürfen die
Variablen $l, E$ und $I_y$ nicht Null sein. Ich denke, dass können wir von denen
auch sagen. QED.
\end{document}
